\scanpage{7}
\mspage{3}
% archive foliation: 6
$\unsure{u}$ \deleted{\unsure{tel}} on a $u^2\in V^*$, donc on a une application
\[
 \operatorname{Circ}(I)\longrightarrow V^*
\]
de degré 2 (qui définit donc une structure de \uline{cube}, avec $V^*$ comme
ens des directions des faces, $\operatorname{Circ}(I)$ comme l'ens des six faces
--- en fait ce cube est canoniquement orienté ; cf plus bas.

\margin{\textcircled{$*$} On peut dire aussi que les six éléments d'ordre 4
se répartissent par paires $\{u,u^{-1}\}$ suivant les trois sous-groupes
\inserted{cycliques} d'ordre 4 \deleted{de} $G=\unsure{\mathcal{G}_I|\mathfrak{S}_I}$,
\unsure{biunivoquement} \deleted{\unsure{en}}\inserted{\unsure{en corresp.}}
avec les éléments de $V^*$ \inserted{(voir au-dessous !)}.}
\margin{$I$ est l'ens des quatre diagonales du cube,}

L'isom. canonique
\[
(\omega)\qquad \underline{\omega}(I)\simeq\underline{\omega}(J)
\]
On a
\[
 \unsure{\mathcal{G}_{I\,\mathrm{ab}}|\mathfrak{S}_{I\,\mathrm{ab}}}
 \xrightarrow{\sim}
 \unsure{\mathcal{G}_{J\,\mathrm{ab}}|\mathfrak{S}_{J\,\mathrm{ab}}}\simeq\{\pm1\}
\]
les isom. étant donnés par la signature, d'où la possibilité de définir un iso.
\deleted{can.}\inserted{fonctoriel} $(\omega)$, dont le choix reste
\unsure{soumis cependant} à convention (deux choix possibles).
Quand on fixe un élément $i\in I$, alors on a
\[
 \unsure{\underline{\omega}_i(I)|\underline{\omega}_1(I)}
 \xrightarrow{\sim}
 \underline{\omega}(\underbrace{I\setminus\{i\}}_{\simeq J})
 \deleted{\unsure{\simeq I\setminus\{i\}}}
\]
en associant à tout ordre \inserted{total} \deleted{\unsure{de}} de $J$
\deleted{\unsure{$J$}} l'ordre \deleted{induit}\inserted{total} de $I$ qui sur $J$ est l'ordre donné,
et qui admet $i$ comme \uline{premier} élément (la convention opposée aurait été
de prendre $i$ comme \uline{dernier} élément).

\margin{$*$
\diagram{\begin{tikzcd}[column sep=large]
V^* \arrow[r,shift left,"\sim","x\mapsto\operatorname{Cent}_G(x)"']
&
\begin{array}{c}\text{Sous-groupes cycliques}\\
\text{d'ordre 4 de }G=\unsure{\mathcal{G}_I|\mathfrak{S}_I}\end{array}
\arrow[l,shift left,"\sim"]
\\
\begin{array}{c}\text{l'unique élément }\ne1\\\text{de }\Gamma\cap V\end{array}
\arrow[r,leftrightarrow]
& \Gamma
\end{tikzcd}}}
\margin{\diagram{\begin{tikzcd}[column sep=small,row sep=large]
& V^* \arrow[rr,leftrightarrow] \arrow[dr] & &
\begin{array}{c}\text{Ss-groupes cycliques}\\
\text{d'ordre 4 de }G=\unsure{\mathcal{G}_I|\mathfrak{S}_I}\end{array}
\arrow[dl,leftrightarrow]\\
\unsure{\alpha|x} \arrow[d] &
\begin{array}{c}
D\mapsto\deleted{\unsure{x}}\\
\text{élément non nul}\\\text{du centre de }D
\end{array}\arrow[u]
&
\begin{array}{c}
\text{Ss-groupes }\deleted{\unsure{de}}\text{ de}\\
\text{Sylow de }G=\unsure{\mathcal{G}_I|\mathfrak{S}_I}\\ (D)
\end{array}
\\
\operatorname{Norm}
\end{tikzcd}}}
